% =============================================================================
%  lec25.tex —— 第 25 讲 Lurie 张量积与 smash 积（底本 PDF 第 186–192 页）
% =============================================================================
\lecture{25}{Lurie 张量积与 smash 积}
\begin{center}\small\color{black!55}（底本 PDF 第 186–192 页）\end{center}

本讲我们定义并研究若干个对称幺半 $\infty$-范畴。

\section{笛卡尔对称幺半结构}

\nitem\label{L25-1.1} 设 $\cC$ 是承认有限乘积的 $\infty$-范畴。本节我们给 $\cC$ 赋予一个对称幺半结构 $\cC^{\times}\to\mathrm{Comm}^{\otimes}$，使得乘法函子由
\[
\cC\times\cC\to\cC,\quad (X,Y)\mapsto X\times Y
\]
给出。

\begin{definition}\label{L25-1.2}
设 $\cC^{\otimes}\to\mathrm{Comm}^{\otimes}$ 是 $\cC$ 上的对称幺半结构。若它满足下述性质，则称它是\textbf{笛卡尔的}：
\begin{enumerate}[label=(\arabic*),leftmargin=2.4em,itemsep=0.1em]
\item 单位对象 $\mathbf{1}_{\cC}\in\cC$ 是终对象。
\item 对对象 $X,Y\in\cC$，态射
\[
X\xleftarrow{\ \simeq\ }X\otimes\mathbf{1}_{\cC}\leftarrow X\otimes Y\to\mathbf{1}_{\cC}\otimes Y\xrightarrow{\ \simeq\ }Y
\]
使 $X\otimes Y$ 成为 $X$ 与 $Y$ 在 $\cC$ 中的乘积。
\end{enumerate}
\end{definition}

\begin{theorem}\label{L25-1.3}（HA.2.4.1.19）
函子 $\cC^{\otimes}\mapsto\cC$ 定义了如下两者之间的等价：
\begin{enumerate}[label=(\roman*),leftmargin=2.4em,itemsep=0.1em]
\item （小）笛卡尔对称幺半 $\infty$-范畴构成的 $\infty$-范畴，以及它们之间的对称幺半函子；
\item （小）承认有限乘积的 $\infty$-范畴构成的 $\infty$-范畴，以及它们之间保持有限乘积的函子。
\end{enumerate}
\end{theorem}

\nitem\label{L25-1.4} 特别地，对承认有限乘积的 $\infty$-范畴 $\cC$，其上的笛卡尔对称幺半结构本质唯一，我们把它记为 $\cC^{\times}\to\mathrm{Fin}_{*}$。注意在使用这个记号时，我们还需要给它配备一个等同 $\cC^{\times}_{\langle1\rangle}\simeq\cC$。

\nitem\label{L25-1.5} 我们可以给 $\Cat_{\infty}$ 配备笛卡尔对称幺半结构 $\Cat^{\times}_{\infty}\to\mathrm{Fin}_{*}$，并考虑其中的 $\OO$-代数对象。注意这样的对象先验地不同于 $\OO$-幺半 $\infty$-范畴：前者是函子 $\OO^{\otimes}\to\Cat^{\times}_{\infty}$，而后者是函子 $\OO^{\otimes}\to\Cat_{\infty}$（经由拉直）。为了比较这两个概念，我们需要一个函子
\[
\Cat^{\times}_{\infty}\to\Cat_{\infty},
\]
它把对象 $(D_{1},\dots,D_{n})\in(\Cat^{\times}_{\infty})_{\langle n\rangle}$ 送到 $\prod_{i=1}^{n}D_{i}\in\Cat_{\infty}$。作为合理性检验，注意惰性态射 $(D_{1},\dots,D_{n})\to D_{k}$ 可以被送到投影函子 $D_{1}\times\cdots\times D_{n}\to D_{k}$。结果表明，这可以推广到任意笛卡尔对称幺半结构 $\cC^{\times}$，包括 $\Cat^{\times}_{\infty}$。

\begin{exercise}\label{L25-1.6}
请说服自己：对一般的对称幺半 $\infty$-范畴 $\cC$，不存在类似的函子 $\cC^{\otimes}\to\cC$。
\end{exercise}

\nitem\label{L25-1.7} 为给出所需函子 $\cC^{\times}\to\cC$ 的精确构造，我们需要一些术语。

\begin{definition}\label{L25-1.8}
设 $p\colon\OO^{\otimes}\to\mathrm{Fin}_{*}$ 是 $\infty$-操作子，$\cD$ 是承认有限乘积的 $\infty$-范畴。设
\[
\Mon_{\OO}(\cD)\subseteq\Fun(\OO^{\otimes},\cD)
\]
是由满足下述条件的函子 $M\colon\OO^{\otimes}\to\cD$ 构成的全子 $\infty$-范畴：
\begin{enumerate}[label=(\roman*),leftmargin=2.4em,itemsep=0.1em]
\item 对 $\OO^{\otimes}_{\langle n\rangle}$ 中任意多色对象 $X=(X_{1},\dots,X_{n})$，典范态射 $M(X)\to M(X_{i})$ 诱导一个等价
\[
M(X)\xrightarrow{\ \simeq\ }\prod_{i=1}^{n}M(X_{i}).
\]
\end{enumerate}
$\Mon_{\OO}(\cC)$ 中的对象称为 $\cC$ 中的 $\OO$-幺半群。
\end{definition}

\begin{example}\label{L25-1.9}
当 $\cC=\Cat_{\infty}$ 时，经由拉直，$\Cat_{\infty}$ 中的 $\OO$-幺半群恰好就是 $\OO$-幺半 $\infty$-范畴。
\end{example}

\begin{definition}\label{L25-1.10}
设 $p\colon\OO^{\otimes}\to\mathrm{Fin}_{*}$ 是对称幺半 $\infty$-范畴，$\cD$ 是承认有限乘积的 $\infty$-范畴。设
\[
\Fun^{\times}(\OO^{\otimes},\cD)\subseteq\Mon_{\OO}(\cD)
\]
是由满足下述条件的 $\OO$-幺半群 $M\colon\OO^{\otimes}\to\cD$ 构成的全子 $\infty$-范畴：
\begin{enumerate}[label=(\roman*),start=2,leftmargin=2.4em,itemsep=0.1em]
\item 对位于唯一活动态射 $\langle n\rangle\to\langle1\rangle$ 上方的任意 $p$-余笛卡尔态射 $X=(X_{1},\dots,X_{n})\to Y$，$M$ 下的像是一个等价
\[
M(X)\xrightarrow{\ \simeq\ }M(Y).
\]
\end{enumerate}
\end{definition}

\begin{remark}\label{L25-1.11}
非正式地说，(i) 与 (ii) 说的是：$M$ 把 $\bigotimes_{i=1}^{n}X_{i}$ 送到 $\prod_{i=1}^{n}M(X_{i})$。
\end{remark}

\begin{proposition}\label{L25-1.12}（HA.2.4.1.6）
设 $\cC$ 与 $\cD$ 是承认有限乘积的 $\infty$-范畴。则沿 $\cC\to\cC^{\times}$ 的限制定义了一个等价
\[
\Fun^{\times}(\cC^{\times},\cD)\xrightarrow{\ \simeq\ }\Fun^{\times}(\cC,\cD),
\]
其中
\[
\Fun^{\times}(\cC,\cD)\subset\Fun(\cC,\cD)
\]
是由保持有限乘积的函子 $\cC\to\cD$ 构成的全子 $\infty$-范畴。
\end{proposition}

\nitem\label{L25-1.13} 由此可知，在 $\Fun^{\times}(\cC^{\times},\cC)$ 中存在本质唯一的函子
\[
\pi\colon\cC^{\times}\to\cC,
\]
使得复合 $\cC\to\cC^{\times}\xrightarrow{\ \pi\ }\cC$ 是恒等函子。

\begin{proposition}\label{L25-1.14}（HA.2.4.1.7）
设 $\cC$ 是承认有限乘积的 $\infty$-范畴。我们给 $\cC$ 配备笛卡尔对称幺半结构。则：
\begin{enumerate}[label=(\arabic*),leftmargin=2.4em,itemsep=0.1em]
\item 对任意 $\infty$-操作子 $\OO^{\otimes}\to\mathrm{Fin}_{*}$，与 $\pi\colon\cC^{\times}\to\cC$ 作复合诱导了一个等价
\[
\Alg_{\OO}(\cC)\xrightarrow{\ \simeq\ }\Mon_{\OO}(\cC).
\]
\item 对任意对称幺半 $\infty$-范畴 $\OO^{\otimes}\to\mathrm{Fin}_{*}$，与 $\pi\colon\cC^{\times}\to\cC$ 作复合诱导了一个等价
\[
\Fun^{\otimes}(\OO,\cC)\xrightarrow{\ \simeq\ }\Fun^{\times}(\OO^{\otimes},\cC).
\]
\end{enumerate}
\end{proposition}

\begin{remark}\label{L25-1.15}
注意 (2) 蕴涵定理~\ref{L25-1.3} 中的函子是全忠实的。事实上，在 HA 中，Lurie 先给出了笛卡尔对称幺半结构 $\cC^{\times}$ 的显式构造，再验证它满足命题~\ref{L25-1.12} 与~\ref{L25-1.14} 中的断言，最后由这些断言推出定理~\ref{L25-1.3}。细节见 HA.2.4.1。
\end{remark}

\begin{remark}\label{L25-1.16}
在 HA 中，Lurie 使用了如下术语：
\begin{itemize}[leftmargin=2.2em,itemsep=0.1em]
\item $\Mon_{\OO}(\cC)$ 中的对象也称为 $\OO^{\otimes}$ 上的 \textbf{lax 笛卡尔结构}；
\item $\Fun^{\times}(\OO^{\otimes},\cD)$ 中的对象称为 $\OO^{\otimes}$ 上的\textbf{弱笛卡尔结构}；
\item $\OO^{\otimes}$ 上的\textbf{笛卡尔结构}是指使得复合 $\OO\to\OO^{\otimes}\to\cD$ 为等价的弱笛卡尔结构。
\end{itemize}
\end{remark}

\section{余笛卡尔对称幺半结构}

\nitem\label{L25-2.1} 利用自同构 $\Cat_{\infty}\to\Cat_{\infty}$、$\cC\mapsto\cC^{\op}$，我们得到定义~\ref{L25-1.2} 与定理~\ref{L25-1.3} 的对偶版本：

\begin{definition}\label{L25-2.2}
设 $\cC^{\otimes}\to\mathrm{Fin}_{*}$ 是 $\cC$ 上的对称幺半结构。若它满足下述性质，则称它是\textbf{余笛卡尔的}：
\begin{enumerate}[label=(\arabic*),leftmargin=2.4em,itemsep=0.1em]
\item 单位对象 $\mathbf{1}_{\cC}\in\cC$ 是始对象。
\item 对对象 $X,Y\in\cC$，态射
\[
X\xrightarrow{\ \simeq\ }X\otimes\mathbf{1}_{\cC}\to X\otimes Y\leftarrow\mathbf{1}_{\cC}\otimes Y\xleftarrow{\ \simeq\ }Y
\]
使 $X\otimes Y$ 成为 $X$ 与 $Y$ 在 $\cC$ 中的余乘积。
\end{enumerate}
\end{definition}

\begin{corollary}\label{L25-2.3}
函子 $\cC^{\otimes}\mapsto\cC$ 定义了如下两者之间的等价：
\begin{enumerate}[label=(\roman*),leftmargin=2.4em,itemsep=0.1em]
\item （小）余笛卡尔对称幺半 $\infty$-范畴构成的 $\infty$-范畴，以及它们之间的对称幺半函子；
\item （小）承认有限余乘积的 $\infty$-范畴构成的 $\infty$-范畴，以及它们之间保持有限余乘积的函子。
\end{enumerate}
\end{corollary}

\nitem\label{L25-2.4} 对承认有限余乘积的 $\infty$-范畴 $\cC$，我们把它上面的余笛卡尔对称幺半结构记为 $\cC^{\sqcup}\to\mathrm{Fin}_{*}$。

\nitem\label{L25-2.5} 但要注意，由练习~\ref{L25-1.6}，命题~\ref{L25-1.12} 与~\ref{L25-1.14} 在余笛卡尔情形下不再成立。作为补偿，我们有下述结果：

\begin{proposition}\label{L25-2.6}（HA.2.4.3.9）
设 $\cC$ 是承认有限余乘积的 $\infty$-范畴。我们给 $\cC$ 配备余笛卡尔对称幺半结构。则对任意带单位的 $\infty$-操作子 $\OO^{\otimes}$，遗忘函子
\[
\Alg_{\OO}(\cC)\to\Fun(\OO,\cC)
\]
是等价。
\end{proposition}

\begin{remark}\label{L25-2.7}
设 $A\colon\OO^{\otimes}\to\cC^{\sqcup}$ 是 $\OO$-代数。上述命题说明，函子 $A$ 由它在 $\OO$ 上的限制决定。非正式地说，对 $\OO^{\otimes}$ 中任意活动态射 $X\defeq(X_{1},\dots,X_{n})\to Y$，其像 $A(X)\to A(Y)$ 可以如下描述。首先，这个态射本质唯一地分解为
\[
A(X)\xrightarrow{\ f\ }\bigsqcup_{i=1}^{n}A(X_{i})\xrightarrow{\ g\ }A(Y),
\]
使得 $f$ 在 $\mathrm{Fin}_{*}$ 上是余笛卡尔的，而 $g$ 落在 $\cC$ 中。态射 $g$ 由态射 $A(X_{i})\to A(Y)$ 决定，而按定义，后者必须是 $X_{i}\to(X_{1},\dots,X_{n})\to Y$ 在 $A$ 下的像。这里态射 $X_{i}\to(X_{1},\dots,X_{n})$ 由 $\OO^{\otimes}$ 带单位这一假设保证。
\end{remark}

\begin{exercise}\label{L25-2.8}
描述 $\Fun^{\otimes}(\OO,\cC)\subseteq\Alg_{\OO}(\cC)$ 在 $\Fun(\OO,\cC)$ 中的本质像。
\end{exercise}

\begin{remark}\label{L25-2.9}
事实上，对任意不必承认有限余乘积的 $\infty$-范畴 $\OO$，都存在典范的带单位 $\infty$-操作子 $\OO^{\sqcup}$，使得命题~\ref{L25-2.6} 的类似结论成立。这样的 $\infty$-操作子称为\textbf{余笛卡尔 $\infty$-操作子}。更多细节见 HA.2.4.3。
\end{remark}

\section{$\infty$-范畴的张量积}

\begin{theoremdef}\label{L25-3.1}（HA.4.8.1.15）
设 $\widehat{\Cat}_{\infty}$ 配备笛卡尔对称幺半结构。考虑 $1$-全子 $\infty$-范畴
\[
(\mathrm{Pr}^{L})^{\otimes}\subseteq(\widehat{\Cat}_{\infty})^{\times},
\]
其定义如下：
\begin{itemize}[leftmargin=2.2em,itemsep=0.1em]
\item 对象 $(C_{1},\dots,C_{n})$ 含于 $(\mathrm{Pr}^{L})^{\otimes}$，当且仅当每个 $C_{k}$ 可表现。
\item 定义在 $\alpha\colon\langle m\rangle\to\langle n\rangle$ 上方的态射 $(C_{1},\dots,C_{m})\to(D_{1},\dots,D_{n})$ 含于 $(\mathrm{Pr}^{L})^{\otimes}$，当且仅当对每个 $j\in\langle n\rangle^{\circ}$，相应的函子
\[
\prod_{i\in\alpha^{-1}(j)}C_{i}\to D_{j}
\]
在每一个因子上保持余极限。
\end{itemize}
则 $(\mathrm{Pr}^{L})^{\otimes}\to\mathrm{Fin}_{*}$ 定义了 $\mathrm{Pr}^{L}$ 上的一个对称幺半结构。我们把乘法函子
\[
-\otimes-\colon\mathrm{Pr}^{L}\times\mathrm{Pr}^{L}\to\mathrm{Pr}^{L}
\]
称为 \textbf{Lurie 张量积}。
\end{theoremdef}

\nitem\label{L25-3.2} 按定义，对 $C_{i}\in\mathrm{Pr}^{L}$，$(\mathrm{Pr}^{L})^{\otimes}\to\mathrm{Fin}_{*}$ 中存在余笛卡尔态射 $(C_{1},\dots,C_{n})\to\bigotimes_{i=1}^{n}C_{i}$。这个态射在更大的纤维化 $(\widehat{\Cat}_{\infty})^{\times}\to\mathrm{Fin}_{*}$ 中几乎绝不是余笛卡尔的。由此可知，在 $(\widehat{\Cat}_{\infty})^{\times}$ 中存在本质唯一的分解
\[
(C_{1},\dots,C_{n})\to\prod_{i=1}^{n}C_{i}\to\bigotimes_{i=1}^{n}C_{i},
\]
使得第一个态射是余笛卡尔的，而第二个落在 $\widehat{\Cat}_{\infty}$ 中。

\begin{exercise}\label{L25-3.3}
对任意 $D\in\mathrm{Pr}^{L}$，以 $\prod_{i=1}^{n}C_{i}\to\bigotimes_{i=1}^{n}C_{i}$ 作前复合，诱导了如下两类对象构成的空间之间的等价：
\begin{itemize}[leftmargin=2.2em,itemsep=0.1em]
\item 保持余极限的函子 $\bigotimes_{i=1}^{n}C_{i}\to D$；
\item 在每个因子上保持余极限的函子 $\prod_{i=1}^{n}C_{i}\to D$。
\end{itemize}
若考虑这些函子构成的 $\infty$-范畴，同样的结论也成立。
\end{exercise}

\begin{exercise}\label{L25-3.4}
证明 $\Spc\in\mathrm{Pr}^{L}$ 是单位对象。函子 $\Spc\times\Spc\to\Spc\otimes\Spc\simeq\Spc$ 是什么？
\end{exercise}

\begin{exercise}\label{L25-3.5}
证明：对任意 $C,D,E\in\mathrm{Pr}^{L}$，
\[
\mathrm{LFun}(C\otimes D,E)\simeq\mathrm{LFun}(C,\mathrm{LFun}(D,E)).
\]
特别地，$\mathrm{LFun}(-,-)$ 计算了关于 Lurie 张量积的内蕴映射对象。
\end{exercise}

\begin{exercise}\label{L25-3.6}
对小 $\infty$-范畴 $C_{0}$ 与 $D_{0}$，证明
\[
\mathrm{PShv}(C_{0})\otimes\mathrm{PShv}(D_{0})\simeq\mathrm{PShv}(C_{0}\times D_{0}).
\]
\end{exercise}

\begin{theorem}\label{L25-3.7}（HA.4.8.1.17）
对 $C,D\in\mathrm{Pr}^{L}$，存在典范等价
\[
C\otimes D\simeq\mathrm{RFun}(C^{\op},D),
\]
其中 $\mathrm{RFun}(-,-)$ 表示由保持小极限的函子构成的 $\infty$-范畴。
\end{theorem}

\begin{remark}\label{L25-3.8}
粗略地说，上述等价是如下得到的。对任意检验对象 $E\in\mathrm{Pr}^{L}$，我们有一族随 $E$ 函子变化的函子
\[
\begin{array}{r@{\ }l}
 & \mathrm{LFun}(C\otimes D,E)\\
\simeq & \mathrm{LFun}(C,\mathrm{LFun}(D,E))\\
\to & \mathrm{LFun}(C,\mathrm{RFun}(E,D)^{\op})\\
\simeq & \mathrm{LFun}(C,\mathrm{LFun}(E^{\op},D^{\op}))\\
\simeq & \mathrm{LFun}(E^{\op},\mathrm{LFun}(C,D^{\op}))\\
\simeq & \mathrm{RFun}(E,\mathrm{RFun}(C^{\op},D))^{\op}\\
\leftarrow & \mathrm{LFun}(\mathrm{RFun}(C^{\op},D),E),
\end{array}
\]
其中不可逆的函子是通过取右伴随得到的：
\[
\mathrm{LFun}(X,Y)\to\mathrm{RFun}(Y,X)^{\op},\qquad F\mapsto F^{R}.
\]
这类函子总是全忠实的，且其本质像由可及右伴随函子 $Y\to X$ 组成。细致的考察表明，$\mathrm{LFun}(C\otimes D,E)$ 与 $\mathrm{LFun}(\mathrm{RFun}(C^{\op},D),E)$ 在这些嵌入下的本质像重合。
\end{remark}

\begin{warning}\label{L25-3.9}
等价 $C\otimes D\simeq\mathrm{RFun}(C^{\op},D)$ 并不以朴素的方式关于 $C$ 与 $D$ 具有函子性。
\end{warning}

\begin{exercise}\label{L25-3.10}
设 $F\colon C_{1}\to C_{2}$ 是 $\mathrm{Pr}^{L}$ 中的态射。证明下图交换：
\[
\begin{tikzcd}[column sep=4.4em,row sep=3.0em]
C_{2}\otimes D \arrow[r,"\simeq"] \arrow[d,"(F\otimes\id_{D})^{R}"'] & \mathrm{RFun}(C_{2}^{\op},D) \arrow[d,"{-\circ F^{\op}}"]\\
C_{1}\otimes D \arrow[r,"\simeq"] & \mathrm{RFun}(C_{2}^{\op},D).
\end{tikzcd}
\]
\end{exercise}

\yizhu{原文此图右下角印作 $\mathrm{RFun}(C_{2}^{\op},D)$；但右侧竖直箭头 $-\circ F^{\op}$ 的靶应为 $\mathrm{RFun}(C_{1}^{\op},D)$，否则图表不交换。译文照录原样。\textbf{[来源：I]}}

\begin{exercise}\label{L25-3.11}
设 $G\colon D_{1}\to D_{2}$ 是 $\mathrm{Pr}^{L}$ 中的态射。证明下图交换：
\[
\begin{tikzcd}[column sep=4.4em,row sep=3.0em]
C\otimes D_{2} \arrow[r,"\simeq"] \arrow[d,"(\id_{C}\otimes G)^{R}"'] & \mathrm{RFun}(C^{\op},D_{2}) \arrow[d,"{G^{R}\circ-}"]\\
C\otimes D_{1} \arrow[r,"\simeq"] & \mathrm{RFun}(C^{\op},D_{1}).
\end{tikzcd}
\]
\end{exercise}

\begin{exercise}\label{L25-3.12}
用练习~\ref{L25-3.10} 或~\ref{L25-3.11} 推出：Lurie 张量积 $-\otimes-\colon\mathrm{Pr}^{L}\times\mathrm{Pr}^{L}\to\mathrm{Pr}^{L}$ 在每一个因子上保持余极限。
\end{exercise}

\begin{exercise}\label{L25-3.13}
设 $C,D\in\mathrm{Pr}^{L}$。回忆我们有伴随 $\cD\rightleftarrows\cD_{*}$，其中左伴随把 $X$ 送到 $X_{+}\defeq X\sqcup\{*\}$。
\begin{enumerate}[label=(\arabic*),leftmargin=2.4em,itemsep=0.1em]
\item 证明存在与到 $\mathrm{RFun}(C^{\op},D)$ 的遗忘函子相容的典范等价
\[
\mathrm{RFun}(C^{\op},D_{*})\to\mathrm{RFun}(C^{\op},D)_{*}.
\]
\item 推出存在与从 $C\otimes D$ 出发的函子相容的典范等价 $C\otimes D_{*}\simeq(C\otimes D)_{*}$。特别地，$\Spc_{*}\otimes\Spc_{*}\simeq\Spc_{*}$。
\end{enumerate}
\end{exercise}

\begin{exercise}\label{L25-3.14}
用上一练习证明：对任意 $C\in\mathrm{Pr}^{L}$，
\[
C\otimes\Sptr\simeq\Sptr(C_{*}).
\]
特别地，$\Sptr\otimes\Sptr\simeq\Sptr$。
\end{exercise}

\section{smash 积}

\begin{definition}\label{L25-4.1}
设 $\cC^{\otimes}\to\mathrm{Fin}_{*}$ 是对称幺半 $\infty$-范畴。若交换代数 $A\in\Alg_{\mathrm{Comm}}(\cC)$ 满足 $A\otimes A\to A$ 是等价，则称 $A$ 是\textbf{幂等的}。令
\[
\Alg^{\mathrm{idem}}_{\mathrm{Comm}}(\cC)\subseteq\Alg_{\mathrm{Comm}}(\cC)
\]
为 $\cC$ 中幂等交换代数构成的全子 $\infty$-范畴。
\end{definition}

\begin{proposition}\label{L25-4.2}（HA.4.8.2.9）
设 $\cC^{\otimes}\to\mathrm{Fin}_{*}$ 是对称幺半 $\infty$-范畴。则遗忘函子
\[
\Alg^{\mathrm{idem}}_{\mathrm{Comm}}(\cC)\to\Alg_{\mathrm{Comm}}(\cC)\simeq\Alg_{\mathrm{Comm}}(\cC)_{\mathbf{1}_{\cC}/}\to\cC_{\mathbf{1}_{\cC}/}
\]
是全忠实的，且其本质像包含那些使得
\[
X\simeq\mathbf{1}_{\cC}\otimes X\to X\otimes X
\]
为同构的对象 $\mathbf{1}_{\cC}\to X$。
\end{proposition}

\begin{exercise}\label{L25-4.3}
设 $\cC\defeq\mathrm{Pr}^{L}$ 配备 Lurie 对称幺半结构。证明态射 $\mathbf{1}_{\cC}\simeq\Spc\to\Spc_{*}$ 与 $\mathbf{1}_{\cC}\simeq\Spc\to\Sptr$ 满足命题~\ref{L25-4.2} 中的条件。推出：$\Sptr$ 与 $\Spc_{*}$ 本质唯一地升级为
\[
\Alg_{\mathrm{Comm}}(\widehat{\Cat}_{\infty})\simeq\Mon_{\mathrm{Comm}}(\widehat{\Cat}_{\infty})
\]
中的对象，使得
\begin{itemize}[leftmargin=2.2em,itemsep=0.1em]
\item 这些对称幺半结构与小的余极限相容；
\item 单位对象分别由 $\mathbb{S}\in\Sptr$ 与 $S^{0}\in\Spc_{*}$ 给出。
\end{itemize}
\end{exercise}

\begin{definition}\label{L25-4.4}
我们把上述对称幺半结构的乘法函子记为
\[
-\wedge-\colon\Spc_{*}\times\Spc_{*}\to\Spc_{*},\qquad
-\otimes-\colon\Sptr\times\Sptr\to\Sptr,
\]
并称它们为 \textbf{smash 积}。
\end{definition}

\begin{exercise}\label{L25-4.5}
证明 $\Spc\to\Spc_{*}\to\Sptr$ 是自然的对称幺半函子。
\end{exercise}

\begin{proposition}\label{L25-4.6}（HA.4.8.2.10）
设 $\cC^{\otimes}\to\mathrm{Fin}_{*}$ 是对称幺半 $\infty$-范畴，$A\in\Alg^{\mathrm{idem}}_{\mathrm{Comm}}(\cC)$。则 $\Mod^{\mathrm{Comm}}_{A}(\cC)^{\otimes}$ 是对称幺半 $\infty$-范畴，且遗忘函子
\[
\Mod^{\mathrm{Comm}}_{A}(\cC)^{\otimes}\to\cC^{\otimes}
\]
是全忠实的。此外，对象 $M\in\cC$ 含于本质像中，当且仅当
\[
M\simeq\mathbf{1}_{\cC}\otimes M\to A\otimes M
\]
是同构。
\end{proposition}

\begin{corollary}\label{L25-4.7}
函子
\[
\Mod^{\mathrm{Comm}}_{\Spc_{*}}(\mathrm{Pr}^{L})\to\mathrm{Pr}^{L}
\]
是全忠实的，且其本质像恰由带基点的可表现 $\infty$-范畴组成。此外，这些 $\infty$-范畴对 Lurie 张量积封闭。
\end{corollary}

\begin{corollary}\label{L25-4.8}
函子
\[
\Mod^{\mathrm{Comm}}_{\Sptr}(\mathrm{Pr}^{L})\to\mathrm{Pr}^{L}
\]
是全忠实的，且其本质像恰由稳定可表现 $\infty$-范畴组成。此外，这些 $\infty$-范畴对 Lurie 张量积封闭。
\end{corollary}

\begin{remark}\label{L25-4.9}
注意上述推论说明：$\Sptr$ 是 $\Alg_{\mathrm{Comm}}(\PrSt^{L})$ 中的始对象，这给出了谱的 smash 积的一个内蕴刻画。
\end{remark}

\appendixsection{更多关于 smash 积}

\nitem\label{L25-A.1} $\infty$-范畴 $\Sptr$ 上的 smash 积可以被严格化为某个单纯模型范畴上的对称幺半结构。其中最成功的一个是对称谱的模型范畴。

\nitem\label{L25-A.2} \textbf{建议阅读}\quad [HSS00]。

\subsection*{参考文献}
\par\noindent\hangindent=5.2em\hangafter=1 [HSS00]\hspace{0.6em}Mark Hovey, Brooke Shipley, and Jeff Smith. Symmetric spectra. \emph{Journal of the American Mathematical Society}, 13(1):149--208, 2000.
